
\section{结论}

我们围绕BRCA分子亚型分类任务，在严格重复交叉验证（Repeated Stratified Cross-Validation）评估口径下，系统比较了RNA-only、Meth-only、Concat、MOFA与Stacking五类融合策略，并完成了特征维度敏感性、CV参数敏感性、MOFA因子数敏感性、Stacking变体对比及消融实验共41组实验。该研究的直接目标是回答多组学融合是否在小样本场景下带来稳定、可解释的性能收益。

本研究的主要结论如下。Concat-SVM取得最优综合性能，在本轮TCGA-BRCA数据，样本约150，上，Concat-SVM方法在Accuracy为90.64\%、Balanced Accuracy为90.56\%和Macro-F1为90.03\%三项指标上均表现最优，较RNA-only基线提升约0.86个百分点。简单方法不劣于复杂方法，RNA单模态已达到89.78\%准确率，MOFA达到90.00\%，三者置信区间高度重叠，差距小于1\%。而Stacking所有变体均低于上述基线，最佳Stacking仅86.00\%，提示在样本远小于特征条件下增加融合层级反而引入过拟合风险。Methylation是重要补充而非独立来源，Meth-only准确率仅83.67\%，显著低于RNA-only为89.78\%，但消融实验表明去除Meth后Concat性能下降约8个百分点，从90.64\%降至88.45\%，说明甲基化信息对跨组学互补具有不可替代的贡献。Stacking性能高度依赖模型选择，7种Stacking变体的Macro-F1范围为77.45\%到85.07\%，方差极大。RF+XGB Hybrid取得组内最优。这表明元学习器在小样本高维空间中的泛化能力存疑。防泄露设计是保证可信度的关键，训练折内拟合、测试折独立推断的严格流程确保了结果的可信度。所有特征筛选、标准化、MOFA拟合与Stacking元特征构造均约束在训练折内完成。

本文的贡献主要体现在以下四个方面。研究设计创新，在单癌种BRCA小样本条件下，建立了严格防泄露的统一评估框架。该框架将41组实验配置组织到6大类别，RNA、Meth、Concat、MOFA、Stacking、消融，中，通过共享evaluation配置文件确保评估参数一致性，使不同方法的比较建立在完全公平的基础上。融合链条实证，以可复现方式给出RNA-only到Meth-only到Concat到MOFA到Stacking的递进对比实证，覆盖41组实验共约1500多个测试折。结果表明简单拼接优于复杂融合，Stacking在小样本下存在系统性过拟合，特征维度在500维后趋于饱和。参数敏感性全景分析，首次在同一数据集上系统性地完成了特征维度，100到1500，CV参数，3/5/10折，3到15次重复，MOFA因子数，10到25，Stacking变体，7种base/meta组合，及消融条件，5类，的全景扫描，为方法选择提供了量化依据。错误分析闭环，将类别级误差，尤其LumB易错现象，纳入主讨论，形成面向后续改进的具体问题清单。LumB亚型F1分数最低，约0.855，是后续改进的重点方向。

从理论价值看，本研究为多组学融合研究中的复杂度-性能权衡提供了反直觉的实证证据。已有文献多报道复杂融合方法的优势，但这些结论往往基于大样本或缺乏严格的防泄露评估。本研究通过41组对照实验揭示了在小样本条件下，样本小于200，简单线性融合，Concat+SVM，的稳健性优于复杂的非线性融合，Stacking。这一发现对方法学讨论具有重要参考价值。

从实践意义看，本研究为临床决策支持系统的开发提供了方法论基础。推荐方案，对于BRCA分子亚型分类任务，推荐使用Concat-SVM或MOFA，20factors，作为首选方案，两者准确率均超过90\%，且计算开销远低于Stacking。不推荐方案，在当前样本规模下，不建议使用Stacking作为生产方案，其性能不稳定且训练成本高昂。资源受限场景，若仅有转录组数据可用，RNA-only SVM已可达到接近90\%的准确率，可作为轻量级替代方案。

准确的分子亚型分类直接关系治疗策略选择，Luminal型适用内分泌治疗，HER2+适用靶向治疗，Basal型适用化疗。通过建立稳健的多组学融合框架，本研究为更精确的临床决策支持奠定方法论基础。

本文明确将novelty限定为实证框架创新而非算法发明。在BRCA单癌种任务中实现了可审计的严格防泄露评估流程，通过41组实验构建了完整的融合策略比较基准。首次系统性地报告了Stacking在小样本下的系统性欠拟合现象，挑战了越复杂越好的直觉假设。建立了配置文件驱动的实验管理体系，提升了实验的可复现性和可扩展性。将类别级误差分析与参数敏感性全景图纳入主讨论，形成面向后续改进的具体问题清单。

本研究存在以下局限，需在解读结论时予以考虑。样本量限制，TCGA-BRCA有效样本约150例，统计功效有限，部分方法间差异可能因样本不足而未达到统计显著性。单一数据集，所有实验均在TCGA-BRCA上完成，结论的泛化性需在其他独立队列如METABRIC上验证。分类器选择，主线实验统一使用SVM/XGB，未穷尽所有分类器组合，其他分类器如深度学习的表现未知。HER2类别样本稀少，HER2+样本在过滤后被剔除或数量极少，四分类验证尚不完整。

面向进一步研究，下一步重点包括。外部验证，引入METABRIC或其他BRCA独立队列进行跨平台验证，确认Concat/MOFA优势是否在独立数据上可复现。深度学习方法探索，引入Transformer或Graph Neural Network等架构，探索端到端多组学融合的可能性。生物学解释强化，补充关键特征与KEGG/GO富集分析，建立特征-通路-亚型的文献映射，将性能差异转化为有意义的生物学洞察。LumB分类难题攻关，深入分析LumB易错样本的分子特征与临床表型，探索针对性改进策略如引入拷贝数变异数据。

